\documentclass[10pt]{article}
\usepackage[margin=0.7in]{geometry}
\usepackage{amsmath,amssymb,amsthm,array,booktabs,mathtools,hyperref}
\hypersetup{hidelinks}
\newtheorem{theorem}{Theorem}
\newtheorem{lemma}{Lemma}
\newcommand{\half}{\tfrac12}
\newcommand{\ps}{\succeq}
\newcommand{\D}{\mathcal D}
\newcommand{\E}{\mathcal E}
\newcommand{\W}{\mathcal W}
\newcommand{\bc}{b_{\mathrm c}}
\newcommand{\Zc}{\mathcal Z}
\newcommand{\R}{\mathbb R}
\setlength{\jot}{3pt}
\title{Completing the state quadratic form at a singular contact}
\author{}\date{}
\begin{document}
\maketitle
\vspace{-2em}

A positive decomposition must retain the kernel at an exact zero. For the
contact considered here, one event square is strictly positive at that zero.
The construction below combines it with both independent state normals.
Completing one normal combines its two positive coefficients. The result,
together with the coefficient of the scalar form $\Psi$ defined below, gives
a positive reserve even at $n=0$. Uniform bounds on the remaining matrix
then give fixed testing functions.

Write the source coordinates as
$c=(d,\lambda,n,m,C,A,D,\alpha,\beta)\geq0$, put
$D_u=1+u^2$, $D_v=1+v^2$, and define the canonical Bell slack
$S_c=F_h/4$ by
\begin{align}
 D_uD_vS_c={}&\lambda D_u(W+vX)^2+nD_uD_vX^2
 +m(vW-X+uvY-uZ)^2 \notag\\
 &+CD_v(uW-Y)^2+AD_uD_vY^2
 +(D+d/2)(uW+uvX-Y-vZ)^2 \notag\\
 &+(\alpha+d/2)D_u(vY-Z)^2
 +(\beta+d/2)D_v(X+uZ)^2-\tfrac d2D_uD_vZ^2 .
 \label{eq:slack}
\end{align}
All four state coordinates are independent. For complex states we add the
same real quadratic form on the real and imaginary parts.

\begin{theorem}\label{thm:weak}
Take the native source defined in Appendix~\ref{app:source}, normalized by
$(d,\lambda,\beta)=(1,0,0)$, on the domain
\begin{equation}
 0<H_o\leq2^{-12},\quad 0<z_o,y_o\leq2^{-64}H_o,\quad
 0<r_o<1,\quad 0<t_o<1+z_or_o,\quad n\geq0.
 \label{eq:native-domain}
\end{equation}
There is a fixed finite family of testing functions, given in
Table~\ref{tab:features}, whose source-dependent positive Gram matrices
represent $S_c$ on
\[
 0\leq u\leq2^{-12},\qquad 0\leq v\leq2^{-25},\qquad w\zeta\geq0.
\]
Here $w=(W+uY)/D_u$ and $\zeta=(Z-uX)/D_u$.
The representation is exact at $n=0$. It implies $S_c\geq0$ for every state
and all $|u|\leq2^{-12}$, $|v|\leq2^{-25}$.
Nonnegative additions to any of the eight event coordinates are admitted
by the fixed enlarged encoding in Section~\ref{sec:encoding}.
\end{theorem}

The source parameters are independent of $u,v$; no lower bound on an opening
and no comparison with a testing opening is assumed. At the native contact,
with $B_s=H_o+y_ot_o$,
\begin{equation}
 (u,v)=(\sqrt{z_o},\sqrt{y_o}),\qquad
 (W,X,Y,Z)=\left(\frac{B_s}{\sqrt{z_oy_o}},-r_o\sqrt{z_o},
                    \frac{H_o}{\sqrt{y_o}},1\right),
 \label{eq:contact}
\end{equation}
we have $S_c=0$, whereas
$c_{\rm test}(aW-bY)=y_ot_o/\sqrt{(1+z_o)(1+y_o)}>0$.
Here $(a,b)=(u,1)/\sqrt{D_u}$ and
$(c_{\rm test},e)=(v,1)/\sqrt{D_v}$. Thus retaining
$C[c_{\rm test}(aW-bY)]^2$ separately would lose this contact.

\section{The two normals and their common reserve}
Use the invertible change
\begin{equation}
 t=\frac{X+uZ}{D_u},\quad r=\frac{Y-uW}{D_u},\quad
 w=\frac{W+uY}{D_u},\quad \zeta=\frac{Z-uX}{D_u};
 \quad
 (W,X,Y,Z)=(w-ur,t-u\zeta,r+uw,\zeta+ut).
 \label{eq:rotation}
\end{equation}
The source change in Appendix~\ref{app:source} gives positive $B,T,j,p,H$
satisfying
\begin{equation}
 B<2^{-11},\quad j,T<2^{-63},\quad 0<p<j,\quad 0<H<1,
 \qquad x=jB=z_o,\quad y=TB=y_o.
 \label{eq:box}
\end{equation}
Put
\begin{gather}
 N=t-pvw,\quad M=-r-Hv\zeta,\quad
 \bc=vw-u\zeta/j,\quad q=v\zeta-Tuw,\quad
 R=\zeta-uv\,w/B,\quad R_u=v(u^2-x)w, \label{eq:rows}\\
 \mu=m(1-p),\quad h_0=\mu+\half,\quad
 M_e=\frac{nh_0j^2}{n+h_0},\quad
 \widehat N=\sqrt{n+h_0}\left(N+\frac{nj}{n+h_0}\bc\right).
 \label{eq:harmonic}
\end{gather}
The positive scalar masses $f_c,q_c,r_c$, and the signed masses $g_c,k_c,d_c$,
are defined explicitly in Appendix~\ref{app:source}.

\begin{lemma}\label{lem:budget}
Let $\Sigma=M_e+f_c/y$. Then
\begin{equation}
 h_0N^2+n(N+j\bc)^2=\widehat N^2+M_e\bc^2,\qquad
 \Sigma>0,\qquad p^2\leq8\Sigma. \label{eq:budget}
\end{equation}
These formulas hold literally at $n=0$.
\end{lemma}
\begin{proof}
The first identity is completion of a square. The source estimates proved
below give, for $\rho=p/j<1$,
\[
 \rho\leq\tfrac52(n+B),\qquad f_c/y\geq j^2B/3.
\]
If $n\leq h_0$, then $M_e\geq nj^2/2$ and
$p^2=j^2\rho^2\leq(5/2)j^2(n+B)
\leq5M_e+(15/2)f_c/y$. If $n\geq h_0$, then
$M_e\geq h_0j^2/2\geq j^2/4$, so $p^2\leq4M_e$.
At $n=0$, $M_e=0$, $\widehat N=\sqrt{h_0}N$ and $\Sigma=f_c/y>0$.
No division by $n$ or $M_e$ occurs.
\end{proof}

To use this reserve, scale only by source quantities:
\begin{gather}
 U=u/\sqrt x,\quad V=v/\sqrt y,\quad \Zc=\zeta/\sqrt{jT},\qquad
 Q_*=Uw-V\Zc,\quad B_*=Vw-U\Zc,\quad R_*=\Zc-UVw, \notag\\
 \Psi=(U-V)^2w^2+\Zc^2(1-UV)^2+2UV(w-\Zc)^2,\qquad
 \E=Q_*^2+R_*^2+V^2(w-UV\Zc)^2.
 \label{eq:circuit}
\end{gather}
Direct expansion gives $\E=(1+V^2)\Psi-V^2B_*^2$.
Writing $\gamma=p^2/\Sigma\leq8$, we obtain
\begin{align}
 \gamma M_ev^4\bc^2+\gamma f_c(yv^2+v^4)\Psi
       -(pv^2\bc)^2&=\gamma f_cyv^2\E, \notag\\
 \gamma TM_eu^2v^2\bc^2+\gamma Tf_cu^2(y+v^2)\Psi
       -(\sqrt T\,puv\bc)^2&=\gamma Tf_cyu^2\E.
 \label{eq:purchases}
\end{align}
Both are polynomial identities in the testing variables; the second is valid
at $v=0$ without dividing the first by $v^2$.
The signed scalar terms combine through
\begin{equation}
 g_cV^2(\Psi-B_*^2)=g_c(\E-\Psi),\qquad
 -\kappa(VQ_*+R_*)^2
 =-2\kappa V^2Q_*^2-2\kappa R_*^2+\kappa(VQ_*-R_*)^2,
 \quad\kappa=-k_c>0. \label{eq:signed}
\end{equation}
These identities retain the contact while redistributing its positive masses.

\section{An exact eleven-row completion}
Set
\[
 a_N=(n+h_0)^{-1/2},\quad
 b_N=\sqrt{\frac{n}{h_0(n+h_0)}},\quad
 \tau=\sqrt{TH(1-H)},\quad
 \rho_Q=T\sqrt{x/q_c},\quad \rho_R=\sqrt{jT/r_c},
\]
and define the column
\begin{equation}
 f=\begin{pmatrix}
 \widehat N\\ M/\tau\\ q/\rho_Q\\ \sqrt{M_e}\bc\\ R/\rho_R\\
 \sqrt{f_c}(U-V)w\\ \sqrt{f_c/(jT)}(1-UV)\zeta\\
 \sqrt{2f_c}(xy)^{-1/4}\sqrt{uv}(w-\Zc)\\
 R_u\\ pv^2\bc\\ \sqrt T\,puv\bc
 \end{pmatrix}.
 \label{eq:f}
\end{equation}
All its coefficients outside testing monomials are source-only.
Let $Q(u,v)=\sum Q_{ab}u^av^b$ be the symmetric matrix specified in
Table~\ref{tab:records}. Each displayed off-diagonal value is a matrix entry,
so its contribution to $f^TQf$ is twice the displayed value.
Define
\begin{align}
 \mathcal R={}&\frac{\mu p}{1-p}N^2+
       \left(\frac{\mu p}{1-p}+h_0+\alpha\right)(uN)^2
       +\half(vN)^2+n[v(N+p\bc)]^2 \notag\\
 &+(C+D+\half)(uM)^2+(C+A+\alpha+\half)(vM)^2+C(uvM)^2
       +d_c(VQ_*)^2+\half p^2y^2(VB_*)^2 \notag\\
 &+g_c[V(w-UV\Zc)]^2+\kappa(VQ_*-R_*)^2.
 \label{eq:retained}
\end{align}
All coefficients in $\mathcal R$ are nonnegative, including
$d_c=H^2(A+C+\alpha+\half)y^2jT+\kappa$.
Put
\begin{equation}
 \W=4\Psi+4U^2(Q_*^2+R_*^2)-V^2(U^2-1)^2w^2.
 \label{eq:W19}
\end{equation}
The explicit positive Gram representation of $\W$ is in
Appendix~\ref{app:parent}. The full completion is
\begin{equation}
 D_vS_c=f^TQf+\mathcal R+16x^2y\W
                  +8\gamma f_cyv^2\E+\gamma Tf_cyu^2\E.
 \label{eq:complete}
\end{equation}
One verifies this equality by substituting Table~\ref{tab:records} and
\eqref{eq:budget}--\eqref{eq:signed} into the complete two-normal expansion
\eqref{eq:fullparent}. Thus it retains both $N$ and $M$ and every original
state direction.

\begin{table}[t]
\centering
\caption{All 44 upper-triangular records of $Q$. Unlisted entries vanish.
Within a row each pair $(i,j)$ has the indicated value.
Write $k_R=K_R+\alpha/B$, with $K_R,K_c$ as in \eqref{eq:normalcoeff}.}
\label{tab:records}
\begin{tabular}{@{}cl@{}}
\toprule
Exponent & Entries of $Q_{ab}$\\ \midrule
$(0,0)$ diagonal&
$1,\ \tau^2(C+D+\half+A),\ 1+g_c/q_c,\ 1,\
1+(g_c-2\kappa)/r_c,$\\
&$1-g_c/f_c-64x^2y/f_c$ (rows 6,7,8), $16,\ 8,\ 1$\\
$(0,0)$ cross&
$(1,9):a_Nk_R,\ (4,9):-b_Nk_R,\ (1,10):a_N/2,\ (4,10):-b_N/2,$\\
&$(2,3):\tau\rho_QA/T,\quad(2,11):(\mu/p)\sqrt{H(1-H)}$\\
$(1,0)$&$(1,5):a_N\alpha\rho_R,\quad(4,5):-b_N\alpha\rho_R$\\
$(0,1)$&$(2,5):\tau\alpha\rho_R$\\
$(2,0)$&$(3,3):-64xy/q_c,\ (5,5):-64xy/r_c,\quad
(i,i):-\gamma Ty\ (i=6,7,8)$\\
$(0,2)$&$(3,3):-2\kappa/(yq_c),\quad
(2,3):\tau\rho_QH(C+A+\alpha+\half),$\\
&$(i,i):-8\gamma y\ (i=6,7,8)$\\
$(1,1)$&$(1,3):a_N(K_c/j)\rho_Q,\quad(3,4):-b_N(K_c/j)\rho_Q,$\\
&$(1,2):(\alpha+\half)\tau a_N,\quad
(2,4):-(\alpha+\half)\tau b_N$\\
$(2,1)$&$(2,5):-\tau BTCH\rho_R$\\
$(2,2)$&$(2,3):\tau\rho_QHC,\quad(i,i):-\gamma T\ (i=4,6,7,8)$\\
$(0,4)$&$(i,i):-8\gamma\ (i=4,6,7,8)$\\
\bottomrule
\end{tabular}
\end{table}

\section{Uniform bounds and fixed gates}
\begin{lemma}\label{lem:bounds}
On \eqref{eq:native-domain}, the matrix in Table~\ref{tab:records} satisfies
$Q_{00}\ps I/4$ and
\begin{equation}
\begin{array}{c|rrrrrrrr}
(a,b)&(1,0)&(0,1)&(2,0)&(0,2)&(1,1)&(2,1)&(2,2)&(0,4)\\ \hline
\|Q_{ab}\|_{\rm op}\leq r_{ab}&1&1&512&8&16&4&16&128 .
\end{array} \label{eq:bounds}
\end{equation}
The bounds hold for $n=0$ and unbounded $n$.
\end{lemma}
\begin{proof}
The explicit source formulas in Appendix~\ref{app:source} give
\begin{equation}
 \frac13\leq
 \frac AT,\ THC,\ T(D+\half)(1-H),\ \frac{\alpha}{B^2},\
 \frac{\mu}{jp},\ \frac{f_c}{Tx^2},\ \frac{q_c}{Tx},\
 \frac{r_c}{TjB^2(1-p)}
 \leq\frac23. \label{eq:scaled}
\end{equation}
Here a list between two bounds means that every member satisfies both.
The polynomial estimate \eqref{eq:NNrem} gives
$n(j-p)/j\leq1$, $\rho\leq5(n+B)/2$. The signed numerator $\mathsf K$
in \eqref{eq:polys} gives $0<\mathsf K\leq2(p+B)$, whence
\begin{equation}
 \frac{\kappa}{r_c}\leq3T(p+B),\quad
 \frac{\kappa}{q_cy}\leq4(p+B),\quad
 \frac{g_c}{f_c}\leq2y,\quad
 |K_c/j|\leq2,\quad |k_R|\leq1.
 \label{eq:normalbounds}
\end{equation}
For the last two bounds use
$K_c/j=p/(2j)+H(\alpha+\half)-\half-n(j-p)/j$ and
\[
 |k_R|\leq\half+\tfrac43\delta^2+\tfrac83\delta<1,\qquad\delta=2^{-10}.
\]
The normalizations therefore satisfy
\begin{gather}
 a_N^2+b_N^2=1/h_0\leq2,\quad a_N+b_N\leq2,\quad
 \tau\rho_Q\leq T,\quad B\rho_R<2,\quad\alpha\rho_R\leq2B,\notag\\
 |(K_c/j)\rho_Q|\leq4\sqrt T,\quad
 (\alpha+\half)\tau\leq\sqrt T/3,\quad
 TH(C+A+\alpha+\half)<1. \label{eq:normalizationbounds}
\end{gather}
For example $\rho_Q\leq\sqrt{3T}$, $\tau\leq\sqrt T/2$,
and $\alpha+\half<2/3$.
Reading the absolute row sums of Table~\ref{tab:records} gives
\[
\begin{array}{c|l}
u&4B\\ v&B\\
u^2&\max\{192B,\ 192/(1-p),\ 8Ty\}\\
v^2&\max\{1+8(p+B),64y\}\\
uv&8\sqrt T\\u^2v&2\sqrt T/3\\
u^2v^2&\max\{2/3,8T\}\\v^4&64 .
\end{array}
\]
Each is below the corresponding fixed bound in \eqref{eq:bounds}.

For the constant matrix, rows $(1,4,9,10)$ have diagonal $(1,1,16,8)$
and cross columns $(a_N,-b_N)^Tk_R$, $(a_N,-b_N)^T/2$.
After subtracting $I/2$, their Schur correction has norm at most
$2(2/31+1/30)=91/465<1/2$.
Rows $(2,3,11)$ have diagonal at least $(1/3,1,1)$: indeed
\[
 TH(1-H)(C+D+\half)\geq(1-H)/3+H/3=1/3.
\]
Their two cross entries have absolute values at most $T,j/3$.
Since $T+j/3<1/12$, this block is at least $I/4$.
Row 5 has diagonal at least $1-6T(p+B)>7/8$; rows 6--8
have diagonal at least $1-2y-192B>7/8$.
These disjoint blocks prove $Q_{00}\ps I/4$.
\end{proof}

\begin{lemma}\label{lem:gates}
The matrix $Q$ has a positive decomposition using only the fixed gates
$\sqrt{1\pm p_{ab}}$, where
$p_{ab}=(2^{12}u)^a(2^{25}v)^b$ and $(a,b)$ runs through
\eqref{eq:bounds}.
\end{lemma}
\begin{proof}
Set $\epsilon_{ab}=2^{-12a-25b}$,
$\alpha_{ab}=r_{ab}\epsilon_{ab}$, $A_{ab}=\epsilon_{ab}Q_{ab}$.
Then $\alpha_{ab}I\pm A_{ab}\ps0$ and
\begin{equation}
 Q=Q_{00}-\sum\alpha_{ab}I+
 \sum_{a,b}\left\{\frac{\alpha_{ab}I+A_{ab}}2(1+p_{ab})
                  +\frac{\alpha_{ab}I-A_{ab}}2(1-p_{ab})\right\}.
 \label{eq:gateidentity}
\end{equation}
On the fixed cap, $0\leq p_{ab}\leq1$. Moreover
\[
 \sum\alpha_{ab}
 =2^{-12}+2^{-25}+2^{-15}+2^{-47}+2^{-33}+2^{-47}+2^{-70}+2^{-93}
 <2^{-10}.
\]
The constant reserve is therefore at least $255I/1024$. All matrices
in \eqref{eq:gateidentity} depend only on the source.
\end{proof}

\section{Fixed functions and the other state signs}
\begin{table}[t]
\centering
\caption{The ordered polynomial/radical list $\psi$ of length 34.
Divide every entry by $\sqrt{D_v}$ to obtain the base testing functions.
The principal list $I_*=(1,5,10,11,13,15,16,17,18,19,20,22,23,24,27,28)$
has length 16. Include $\sqrt{1\pm p_{ab}}\,\psi_{I_*}/\sqrt{D_v}$
for each of the eight exponent pairs.}
\label{tab:features}
\begin{tabular}{@{}cl@{}}
\toprule
Factor& Multipliers, in order\\ \midrule
$t$&$1,v,u,uv$\\
$r$&$1,v,u,uv$\\
$w$&$1,v,u,v^2,uv,u^2,v^3,uv^2,u^2v$\\
$\zeta$&$1,v,u,v^2,uv,uv^2,u^2v,uv^3,u^2v^2$\\
$\sqrt{uv}\,w,\ \sqrt{uv}\,\zeta$&$1$ for each\\
$\sqrt{uw\zeta}$&$1,v,u,uv$\\
$\sqrt{vw\zeta}$&$1,u$\\ \bottomrule
\end{tabular}
\end{table}

Every row of $f$ lies in the source span of $\psi_{I_*}$:
its eleven supports are respectively
\[
 \{1,10,20\},\{5,19\},\{11,19\},\{10,20\},\{13,18\},
 \{10,11\},\{18,22\},\{27,28\},\{10,17\},\{15,23\},\{16,24\}.
\]
All square factors in $\mathcal R$ lie in the source span of $\psi$.
The same holds for each of the three factors in $\E$ after multiplication
by $u$ or $v$, by \eqref{eq:rows}--\eqref{eq:circuit}.
The seven polynomial roots of $\W$ in Appendix~\ref{app:parent} also lie
there. Its remaining roots are proportional to
$\sqrt{vw\zeta}(1-U)$ and
$\sqrt{uw\zeta}(1-UV,V-UV,U-UV)$, using indices 29--34.
Thus \eqref{eq:complete} and \eqref{eq:gateidentity} yield the fixed
34-dimensional base block and sixteen 16-dimensional gate blocks.

\begin{lemma}\label{lem:sign}
The positive-chart certificate implies the cap conclusion of
Theorem~\ref{thm:weak}, including its literal axes.
\end{lemma}
\begin{proof}
Every source summand in \eqref{eq:slack} is invariant under
\[
 (u,v,W,X,Y,Z)\longmapsto
 (\varepsilon u,\eta v,W,\eta X,\varepsilon Y,\varepsilon\eta Z).
\]
The rotated coordinates transform as
$(t,r,w,\zeta)\mapsto(\eta t,\varepsilon r,w,\varepsilon\eta\zeta)$.
Let $K$ be the $d=1$ core in the numerator of \eqref{eq:slack},
and put $h=1+u^2+v^2$, $\Delta=h^2-u^2v^2>0$,
\[
 A_0=\half\begin{pmatrix}h&-uv\\-uv&h\end{pmatrix},\qquad
 B_0=\half\begin{pmatrix}-u^2v&-uv^2\\uv^2&u^2v\end{pmatrix}.
\]
Direct completion, also printed as SJ5 in the companion
\cite{companion}, gives
\begin{align}
 K/D_u={}&\left[\binom tr+A_0^{-1}B_0\binom w\zeta\right]^TA_0
 \left[\binom tr+A_0^{-1}B_0\binom w\zeta\right] \notag\\
 &+\frac{D_uD_v}{2\Delta}
 \left[u^2v^2(w^2+\zeta^2)-2uvh\,w\zeta\right].
 \label{eq:signcore}
\end{align}
Since $A_0\succ0$, $K<0$ forces $uvw\zeta>0$. The gauges then put
$u,v>0$ and $w\zeta>0$, where \eqref{eq:complete} proves positivity.
If $K\geq0$, the core and all eight event terms are already nonnegative.
At $u=0$, $v=0$, $w=0$ or $\zeta=0$, \eqref{eq:signcore} has
nonnegative remainder; no limiting division by an axis is needed.
\end{proof}

\section{The affine encoding and its scope}\label{sec:encoding}
A fixed enlargement makes all event additions transparent. Set
\[
 H_{11}=(X,Y,Z,vW,uvW,uW,uvX,vZ,u^2vW,u^3W,uZ)
\]
and append to it
$(uX,vX,v^2W,uv^2W,vY,uvY,W,uvZ,uY)$, obtaining $H_{20}$.
Use $H_{20}/\sqrt{D_uD_v}$ as one base block and
$\sqrt{1\pm p_{ab}}H_{11}/\sqrt{D_uD_v}$ for
\[
 (a,b)=(0,1),(0,2),(1,0),(1,1),(1,2),(2,0),(2,1),
 (2,2),(3,0),(3,2),(4,0),(4,2),(5,0).
\]
These extra gate blocks may all be zero for the weak native witness.
The companion note uses them for $H_o\geq2^{-12}$.
The eight event rays have Gram witnesses in the base $H_{20}$ alone:
\begin{equation}
\begin{array}{c|l@{\qquad}c|l}
\lambda&e_{18}+e_{13},\ e_6+e_7&
n&e_1,\ e_{12},\ e_{13},\ e_7\\
m&e_4-e_1+e_{17}-e_{11}&
C&e_6-e_2,\ e_5-e_{16}\\
A&e_2,\ e_{20},\ e_{16},\ e_{17}&
D&e_6+e_7-e_2-e_8\\
\alpha&e_{16}-e_3,\ e_{17}-e_{11}&
\beta&e_1+e_{11},\ e_{13}+e_{19}.
\end{array}\label{eq:eventrays}
\end{equation}
Each table entry specifies roots; their outer products are added to the
corresponding source Gram. This proves the addition claim, including $d=0$.

For clarity the affine equations can be specified without the nonlinear
source dictionary. Let $T_k$ be the nine polynomial coefficients of $c_k$
on the right of \eqref{eq:slack}, after \eqref{eq:rotation}.
For each block let $h_b$ be its undivided list. Its cleared product weight
is $1$ for the $H_{20}$ block, $1\pm p_{ab}$ for an $H_{11}$ gate block,
$D_u$ for the weak base, and $D_u(1\pm p_{ab})$ for a weak gate block.
Equating every coefficient in
\begin{equation}
 \sum_b\omega_b h_b^TG_bh_b=\sum_k c_kT_k,\qquad G_b\ps0,\quad c\geq0
 \label{eq:affine}
\end{equation}
defines a rational affine system. Off-diagonal Gram variables enter with
factor two. Reduce products modulo
$s_u^2=u,\ s_v^2=v,\ s_z^2=w\zeta$, where the three symbols denote the
nonnegative square roots. Their square classes are independent, so this
coefficient comparison is exact on the positive open chart.

One may average the eight formal sign changes of $s_u,s_v,s_z$ in
\eqref{eq:affine}. They fix all targets and weights and preserve PSD by
diagonal conjugation. Cross-character entries vanish. The weak base splits
as $26+2+4+2$ and each weak gate as $14+2$; the other blocks are unchanged.
There are then 63 Gram blocks, total Gram order 596, 4021 Gram variables
and 619 coefficient equations. Nine scalar source signs give order 605.
These are counts of the specified system, not ranks or additional admission
hypotheses. The literal rational arrays and their character reduction
accompany this note.

The same fixed system admits finite limits of its feasible source coefficients.
Indeed every block's feature list is independent on the strict chart.
A common finite collection of chart evaluations therefore gives a positive
definite evaluation Gram in each block. Summing \eqref{eq:affine} at these
points bounds all PSD traces when the source coefficients converge; a
subsequence of the Gram tuples then gives the limiting identity.

Finally the cap contains the complement left by the two preserved regional
certificates of \cite{regional}:
\[
 b>1-2^{-26}\ \Longrightarrow\
 u^2<(1-2^{-26})^{-2}-1<2^{-24},\qquad
 |c_{\rm test}|<2^{-26}\ \Longrightarrow\
 v^2<\frac{2^{-52}}{1-2^{-52}}<2^{-50}.
\]
Those certificates cover the remaining tests, including projective endpoints,
on the same native envelope. Appending their unchanged order-20 and order-10
pencils gives total order 635, 4024 variables and 619 equations.
The joined cone $K_{635}$ is closed. The cap Grams are bounded along every
convergent source sequence by the preceding evaluation argument. Put
$\rho=\tau-d/2$ and $B_0=1-2^{-26}$. The order-20 pencil has diagonal
entries $2\rho/B_0^2$, $q/2$, $p-q$, $d/2-p$ and a principal minor
$\rho(2D-\rho)/B_0^2$. Their nonnegativity gives
$0\leq\rho\leq2D$ and $0\leq q\leq p\leq d/2$, including $\rho=0$.
Passing to a subsequence therefore preserves both regional pencils.
Together with the strong-contact proof, we obtain
\[
 \overline{\mathcal E+\operatorname{cone}(\mathcal S_{\rm native})}
 \subseteq K_{635}\subseteq C_{\rm can},
\]
Here $\mathcal E$ is the positive event cone, $\mathcal S_{\rm native}$
the native envelope with $0<H_o<1$, and $C_{\rm can}$ the canonical
valid-source cone. Whether $K_{635}=C_{\rm can}$ remains open.

\appendix
\section{Native coefficients and uniform source estimates}\label{app:source}
The following dictionary specifies the native source used above; it is the
RSR source change of \cite{companion}. Put
\begin{gather}
 B_s=H_o+y_ot_o,\quad h_s=H_o/B_s,\quad c_s=1+z_or_o,\quad
 R_s=1+h_sz_o,\notag\\
 B=B_sR_s/c_s,\quad j=z_oc_s/(B_sR_s),\quad
 T=y_oc_s/(B_sR_s),\quad
 p=z_o(1-r_o)/(B_sR_s),\quad H=t_o/c_s.
 \label{eq:sourcechange}
\end{gather}
Then $x=jB=z_o$, $y=TB=y_o$, $p/j=(1-r_o)/c_s$ and
\eqref{eq:box} follows directly. Define the polynomials
\begin{align}
 K&=j+1-p,\quad E=1+B(HT+p-1),\quad W_0=(1-H)j+1-p,\notag\\
 P&=j+1+B((j+1)^2-p)
 +TB[Hj+j+1+B(Hj(2j+2-p)+p(2j+1-p))]\notag\\
 &\hspace{2em}+HjT^2B^2(1+Bp+BHj),\notag\\
 L&=K+T(W_0-H)+BT[pK+Hj(1-p)]+BHT^2(W_0-1),\notag\\
 \mathsf N&=p+B(-j^2+3jp-j-p^2+p)+BT(Hj^2+Hj-j^2+jp-j+p)\notag\\
 &\quad+B^2T(2Hj^2p-Hj^2-Hjp^2+2Hjp-2j^2p+3jp^2-jp-p^3+p^2)
       +B^2HT^2j(Hj-j+p)(1+Bp),\notag\\
 F_1&=-2BHjp+BHj+2Bjp-2Bp^2+2Bp-Hj-Hp-H+j+1-p^2/j+p/j,\notag\\
 F_2&=-B^2H^2jp+B^2Hjp-B^2Hp^2+B^2Hp-BH^2j-BH^2p
                 +BHj-BHp^2/j+BHp/j,\notag\\
 F&=K+TF_1+T^2F_2,\notag\\
 \mathsf K&=pW_0+B(1-p)[j(H(p+1)-p)-(1-p)(p-H)]\notag\\
 &\quad+TBH[(p+1)(1-H)j+(1-p)(BHj+p+1-H)]
                    +T^2B^2H^2j(1-H).
 \label{eq:polys}
\end{align}
The six native weights and the scalar masses are
\begin{align}
 \mu&=\frac{jp(1+y)(1+x)(1+Hy)}{2P},&
 m&=\frac{\mu}{1-p},\notag\\
 n&=\frac{\mathsf N}{2(j-p)P},& C&=\frac{L}{2HTP},\notag\\
 A&=\frac{T(1+x)(1+Hy)(K+Hjy)}{2(1-HT)P},\notag\\
 D&=\frac{(1+Hy)(K+Hjy)[1-y((1-H)j-p)]}{2T(1-H)P},&
 \alpha&=\frac{B^2(K+Hjy)L}{2EP},\notag\\
 f_c&=\frac{Tx^2F}{2P},& q_c&=Ax-f_c,\notag\\
 r_c&=\frac{B^2jT(1-p)(1+x)(1+Hy)L}{2EP},\notag\\
 g_c&=\frac{HT^2j^2B^3L}{2P},&
 k_c&=-\frac{T^2B^2j(1+x)\mathsf K}{2EP},\notag\\
 d_c&=H^2(A+C+\alpha+\half)y^2jT-k_c.
 \label{eq:source}
\end{align}
The condition $n\geq0$ is exactly $\mathsf N\geq0$.

Here are explicit finite estimates sufficient for all bounds in the proof.
For a polynomial after $p=js$, let its $\delta$-envelope be the sum of the
absolute coefficients weighted by
$\delta^{\deg_B+\deg_T+\deg_j}$, with $\delta=1/1024$ and $0\leq H,s\leq1$.
The polynomials in \eqref{eq:polys} give
\begin{equation}
\begin{array}{c|ccccc}
 &P-1&L-1&F-1&F_1&F_2/B\\ \hline
\text{envelope}\leq&
\frac{1155174955302256641}{590295810358705651712}&
\frac{2149057539}{549755813888}&
\frac{5633901886178307}{1152921504606846976}&
\frac{1612711427}{536870912}&
\frac{1077937155}{1073741824}
\end{array}\label{eq:envelopes}
\end{equation}
The first three are below $1/128$, the last two below $4,2$.
In particular $1\leq P\leq129/128$, $127/128\leq L,F\leq129/128$,
and $1-\delta\leq E\leq1$.
Put $\ell=127/128$, $h=129/128$,
$a_-=1/(2h)$,
$a_+=(1+\delta^2)^2(1+2\delta)/[2(1-\delta)]$ and
$d_+=(1+\delta^2)(1+2\delta)(1+\delta^3)/2+\delta/2$.
Termwise substitution into \eqref{eq:source} yields
\begin{equation}
\begin{array}{c|cc}
\text{ratio}&\text{lower}&\text{upper}\\ \hline
 A/T&a_-&a_+\\
 THC&\ell/(2h)&h/2\\
 T(D+\half)(1-H)&(1-\delta^3)/(2h)&d_+\\
 \alpha/B^2&\ell/(2h)&(1+2\delta)h/[2(1-\delta)]\\
 \mu/(jp)&a_-&(1+\delta^2)^3/2\\
 f_c/(Tx^2)&\ell/(2h)&h/2\\
 q_c/(Tx)&a_--2\delta^2/3&a_+\\
 r_c/[TjB^2(1-p)]&\ell/(2h)&(1+\delta^2)^2h/[2(1-\delta)]
\end{array}\label{eq:ratiotable}
\end{equation}
Every lower endpoint exceeds $1/3$ and every upper endpoint is below $2/3$.
For example use $x,y\leq\delta^2$,
$1\leq K+Hjy\leq1+2\delta$ and
$1-\delta^3\leq1-y[(1-H)j-p]\leq1+\delta^3$;
for $q_c$ use $q_c/(Tx)=A/T-xf_c/(Tx^2)$.

For the only estimate involving $n$, write
\begin{gather}
 \mathsf N=p[1+B(3j+1-p)+BT(j+1)]
       -Bj(j+1)[1+(1-H)T]+\mathcal R_N,\notag\\
 \frac{\mathcal R_N}{B^2Tj^2}
 =H(2s-1)-s+s^2+j(2Hs-Hs^2-2s+3s^2-s^3)
                 +HT(H-1+s)(1+Bjs),\quad s=p/j.
 \label{eq:NNrem}
\end{gather}
The last expression has absolute value at most
$5+9\delta+\delta(1+\delta^2)<6$.
Thus $\mathsf N\geq p-2Bj(j+1)$ and
$\mathsf N/j\leq1+B(3j+1)+BT(j+1)+6B^2T(j+1)<2$.
These give $n(j-p)/j\leq1$ and $p/j\leq(5/2)(n+B)$.
Finally
\begin{align*}
 \mathsf K-pW_0+Bp(1-p)(j+1)
 ={}&B(1-p)[Hj(p+1)+H(1-p)+p^2]\\
 &+TBH[(p+1)(1-H)j+(1-p)(BHj+p+1-H)]\\
 &+T^2B^2H^2j(1-H)\geq0 .
\end{align*}
It follows that $\mathsf K>0$; its displayed polynomial bounds it above by
$ (1+\delta)p+(1+3\delta+3\delta^2+\delta^3+\delta^4)B
\leq2(p+B)$. Inserting these two bounds into \eqref{eq:source} gives
\eqref{eq:normalbounds}. All source denominators remain strictly positive;
undefined zero-denominator native parameter values are not assigned values.

\section{The full parent and the retained circuit}\label{app:parent}
For a direct verification of \eqref{eq:complete}, set
\begin{align}
 K_c&=p/2+j[H(\alpha+\half)-\half]+n(p-j),&
 K_R&=-\mu-(\alpha+\half)(1-p)+(T/j)K_c,\notag\\
 C_q&=A/T+v^2[CH(1+u^2)+AH+(\alpha+\half)H],\notag\\
 L_N&=nj\bc+npv^2\bc-j\alpha\bc+\half pv^2\bc
                      +K_RR_u+(K_c/j)uvq,\notag\\
 L_M&=C_qq+\mu uv\bc+v(\alpha-BTCHu^2)R,\notag\\
 H_N&=h_0(1+u^2)+\half v^2+\alpha u^2+n(1+v^2),\notag\\
 H_M&=(1+u^2)[C(1+v^2)+D+\half]+A(1+v^2)+(\alpha+\half)v^2.
 \label{eq:normalcoeff}
\end{align}
The scalar part and both full normals are
\begin{align}
 F_s={}&q_cQ_*^2+nj^2yB_*^2+r_cR_*^2+d_cV^2Q_*^2
 +(np^2y^2+\half p^2y^2-g_c)V^2B_*^2\notag\\
 &+f_c\Psi+g_cV^2\Psi+k_c\Zc^2(V^2-1)^2,\notag\\
 D_vS_c={}&F_s+2L_NN+2L_MM+H_NN^2+H_MM^2
       +2(\alpha+\half)uvNM+\frac{\mu p}{1-p}(1+u^2)N^2.
 \label{eq:fullparent}
\end{align}
Substitution of \eqref{eq:source} and \eqref{eq:rotation} in
\eqref{eq:slack} proves \eqref{eq:fullparent} coefficient by coefficient,
equivalently RFN8/RFN12 of \cite{companion}. The cancellations needed for
this substitution, with $K_d=CH-(D+\half)(1-H)$, include
\begin{align*}
 0&=T(K_d+AH+\half)-A,\\
 0&=jK_d+\mu+xTCH,\\
 0&=jK_d+xTCH+T[CH+AH+(\alpha+\half)H]
                -A-(\alpha+\half)(1-p)+\alpha/B,\\
 0&=np-\mu+p/2+[-\mu-(\alpha+\half)(1-p)]x+j(\alpha-n)+yK_c .
\end{align*}
These are rational identities in the printed dictionary; no contact equation
on the state is imposed.

For completeness the RFN19 remainder has the following literal Gram form.
Let
\[
 t_7=\begin{pmatrix}
 \Zc-U^2Vw\\Vw-U^2Vw\\Uw-U^2Vw\\U\Zc-U^2Vw\\
 UVw-U^2Vw\\UV\Zc-U^2Vw\\U^2w-U^2Vw
 \end{pmatrix},\qquad
 t_3=\begin{pmatrix}1-UV\\V-UV\\U-UV\end{pmatrix}.
\]
Then
\begin{equation}
 \W=t_7^TG_7t_7+4Vw\Zc(1-U)^2+Uw\Zc\,t_3^TG_3t_3,
 \label{eq:RFN19}
\end{equation}
where
\[
 G_7=\begin{pmatrix}
4&-2&-3/2&0&-1/2&0&1/2\\
-2&3&0&-3/2&0&-1&2\\
-3/2&0&4&3/2&-2&-3&0\\
0&-3/2&3/2&4&-3/2&0&-3/2\\
-1/2&0&-2&-3/2&3&2&-1\\
0&-1&-3&0&2&8&-4\\
1/2&2&0&-3/2&-1&-4&4
\end{pmatrix},\qquad
 G_3=\begin{pmatrix}3&-2&-2\\-2&2&2\\-2&2&3\end{pmatrix}.
\]
Their leading principal minors are respectively
$(4,8,101/4,1049/16,907/16,3279/16,743/16)$ and $(3,2,2)$.
Thus both matrices are positive definite. Since $U,V\geq0$ and
$w\Zc\geq0$, \eqref{eq:RFN19} is a sum of squares in precisely the
radicals included in Table~\ref{tab:features}.

\begin{thebibliography}{2}
\bibitem{companion}
\emph{Complete proof companion}, \path{latex/certificates/c2_short.tex}:
source identities RSR2--RSR19, full-parent identities RFN4--RFN12 and
RFN18b--RFN19, and sign completion SJ1--SJ7.
\bibitem{regional}
\emph{Fixed square decompositions away from a singular contact},
\href{../fixed-c2-gates/Fixed_C2_Completions.pdf}{regional identities
and native-source estimates}.
\end{thebibliography}
\end{document}
